% Full arbitrary-lattice lifts of the original CCM norm and Weil transfer.
\section{Galois norm and Weil transfer with every support label retained}
\label{sec:cc-galois-transfer-full-lattice}

\paragraph{Human source and exact input.}
Alain Connes, Caterina Consani and Matilde Marcolli,
\href{https://arxiv.org/abs/math/0703392v1}{\emph{The Weil proof and the
geometry of the adeles class space}}, original author file
\nolinkurl{Yuri70v2.tex}, lines 3514--3853, provide the original
construction used here. Its source labels are
\texttt{Normmap}, \texttt{AKGemb}, \texttt{explnorm},
\texttt{CWGext}, \texttt{rulesW}, \texttt{transferhom},
\texttt{taumap}, \texttt{leftrightact}, \texttt{tauequiv},
\texttt{Xitau}, and \texttt{tauXiC1}.
The author TeX SHA256 is
\texttt{b1f694cb934dd4fe\allowbreak{}f829287f1a55dd8d\allowbreak{}e21c67a481e55245\allowbreak{}3c2f7e8fe0ae6fb5}.
The class-field assertion used as an input is the surjectivity of
the specified relative Weil-group transfer. CCM cite Andr\'e Weil,
\emph{Sur la th\'eorie du corps de classes}, J. Math. Soc. Japan
\textbf{3} (1951), 1--35, for that assertion. Their local norm
references are to Weil's \emph{Basic Number Theory}.
No independent rereading or new proof of those class-field theorems
is claimed here. The carrier maps, their exact fibres, the typed
group-transfer formula, and the lifted orbit maps below are proved.
The carriers are the arbitrary-$L$ carriers PLP1--14; the local
product is not substituted for the original common-label object.

\subsection{The original field maps and the two different carriers}

Let $E/K$ be a finite Galois extension of global fields, with
$G=\operatorname{Gal}(E/K)$ and $n=[E:K]$. Denote its map of places
by $\pi:\Sigma_E\to\Sigma_K$. For $w\mid v$ put
$d_w=[E_w:K_v]$. Galois transitivity gives
\begin{equation}\tag{GTL1}
 \sum_{w\mid v}d_w=n,
 \qquad G_w=\operatorname{Gal}(E_w/K_v),\qquad |G_w|=d_w.
\end{equation}
Use the original diagonal inclusion and local norm formula
\begin{equation}\tag{GTL2}
 \begin{aligned}
 \iota:\mathbb A_K&\longrightarrow\mathbb A_E,
                   &\iota(a)_w&=a_{\pi(w)},\\
 N:\mathbb A_E&\longrightarrow\mathbb A_K,
                   &(Nx)_v&=\prod_{w\mid v}N_{E_w/K_v}(x_w).
 \end{aligned}
\end{equation}
The inclusions $K_v\to E_w$ in the first line are the original
field inclusions. Integral local elements have integral norm,
so the second line respects the restricted-product condition.
The Galois action on adeles is
$(\gamma x)_w=\gamma(x_{\gamma^{-1}w})$, with the indicated
isomorphism of local fields. Its fixed subring is $\iota(\mathbb A_K)$.

The local field norm is the determinant of multiplication on the
$d_w$-dimensional $K_v$-vector space $E_w$. It is multiplicative,
equals zero exactly at zero, and sends a scalar $a_v$ to $a_v^{d_w}$.
Regrouping the full product over $G$ by the places above $v$ gives
\begin{equation}\tag{GTL3}
 \iota(Nx)=\prod_{\gamma\in G}\gamma x,
 \qquad N\iota(a)=a^n.
\end{equation}
In this regrouping each place above $v$ occurs $d_w$ times,
with its local automorphisms, and GTL1 supplies the exact exponent.

Let $L$ be the original nontrivial bounded distributive lattice.
For $F=K,E$ define the two specified semirings
\begin{equation}\tag{GTL4}
 \begin{gathered}
 S_{F,c}=G_L(\mathbb A_F),\qquad
 S_{F,\mathrm{loc}}=\prod_{v\in\Sigma_F}'G_L(F_v),\qquad
 \Delta_F(a,\lambda)=((a_v,\lambda))_v,\\
 G_L(R)=\{(0,\lambda):\lambda\in L\}\cup(R\times\{1_L\}),\\
 (a,\lambda)+(b,\mu)=(a+b,\lambda\vee\mu),\quad
 (a,\lambda)(b,\mu)=(ab,\lambda\wedge\mu).
 \end{gathered}
\end{equation}
The finite restricted factors are $G_L(\mathcal O_v)$.
Every zero label belongs to these factors, so the restricted
condition is entirely the original amplitude condition. The
map $\Delta_F$ is the injective unital semiring homomorphism
of PLP3: its image consists of the tuples with one common label.
At a nonzero adele that common label is forced to be top, including
at any of its zero coordinates.

The exact lifted inclusion and norm maps are
\begin{equation}\tag{GTL5}
 \begin{aligned}
 \iota_c(a,\lambda)&=(\iota(a),\lambda),&
 N_c(x,\lambda)&=(Nx,\lambda),\\
 (\iota_{\rm loc}(a,\boldsymbol\mu))_w
       &=(a_{\pi(w)},\mu_{\pi(w)}),&
 (N_{\rm loc}(x,\boldsymbol\lambda))_v
       &=\left(\prod_{w\mid v}N_{E_w/K_v}(x_w),
                         \bigwedge_{w\mid v}\lambda_w\right).
 \end{aligned}
\end{equation}
All meets here are finite, so no completeness of $L$ is assumed.
A nonzero target norm at $v$ forces every $x_w$ above it to be
nonzero and hence every input label to be top. This proves that
$N_{\rm loc}$ lands in the stated constrained carrier. The same
argument proves the common-label constraint for $N_c$.
Both inclusions are injective unital semiring maps. Both norm
maps preserve multiplication, the multiplicative unit, and the
additive-zero element $\tau$; they are multiplicative monoid maps.

On the local carrier the Galois action is
\begin{equation}\tag{GTL6}
 \gamma(x,\boldsymbol\lambda)_w
   =\bigl(\gamma(x_{\gamma^{-1}w}),\lambda_{\gamma^{-1}w}\bigr).
\end{equation}
It is a semiring automorphism. On the common carrier it preserves
the one label. The fixed local tuples have constant labels over
each fibre of $\pi$ and fixed amplitudes, so they are exactly the
image of $\iota_{\rm loc}$; the common fixed points are exactly
the image of $\iota_c$. Finite idempotence of meet, including its
$d_w$ repeated occurrences, now gives all the exact squares
\begin{equation}\tag{GTL7}
 \begin{gathered}
 N_{\rm loc}\Delta_E=\Delta_KN_c,\qquad
 \iota_{\rm loc}\Delta_K=\Delta_E\iota_c,\\
 N_c\iota_c(a,\lambda)=(a^n,\lambda)=(a,\lambda)^n,\\
 N_{\rm loc}\iota_{\rm loc}(a,\boldsymbol\mu)
                         =(a^n,\boldsymbol\mu),\\
 \iota_cN_c(X)=\prod_{\gamma\in G}\gamma X,\qquad
 \iota_{\rm loc}N_{\rm loc}(X)=\prod_{\gamma\in G}\gamma X.
 \end{gathered}
\end{equation}
Thus the norm/diagonal and inclusion/diagonal squares commute on
the full original carriers, not merely on their amplitudes.

\subsection{Exact zero-pattern fibres and the additive comparison}

For a fixed adele $x\in\mathbb A_E$ let $Z(x)=\{w:x_w=0\}$.
Field norm detects zero, so
\begin{equation}\tag{GTL8}
 Z(Nx)=\pi(Z(x)).
\end{equation}
The local amplitude fibre over $x$ is $L^{Z(x)}$.
For each $v\in\pi(Z(x))$ set
$A_v=Z(x)\cap\pi^{-1}(v)$, a nonempty finite set. The exact
map of these amplitude fibres and its fibres at an output
$\boldsymbol\mu\in L^{\pi(Z(x))}$ are
\begin{equation}\tag{GTL9}
 \begin{gathered}
 L^{Z(x)}\longrightarrow L^{\pi(Z(x))},\qquad
 (\lambda_w)_w\longmapsto
                 (\bigwedge_{w\in A_v}\lambda_w)_v,\\
 \prod_{v\in\pi(Z(x))}
 \left\{(\lambda_w)_{w\in A_v}:
                         \bigwedge_{w\in A_v}\lambda_w=\mu_v\right\}.
 \end{gathered}
\end{equation}
Top labels at the nonzero coordinates have been inserted, not
deleted from the carrier. This fibre map is surjective: select one
zero place above each such $v$, put $\mu_v$ there and put top at
the remaining zero places. This choice is only a set-theoretic
section of the stated finite-fibre meet maps, not a new canonical
Galois choice.

For the common carrier, a fixed nonzero $x$ admits only its top
label, even when $Nx=0$. Such an element maps to $e=(0,1_L)$
when its entire norm adele vanishes. At the all-zero adele the
map is $(0,\lambda)\mapsto(0,\lambda)$. Consequently the
complete inverse fibre of $(0,\mu)$ for $\mu<1_L$ is the
single point $(0,\mu)$; the fibre of $e$ is
$\{(x,1_L):Nx=0\}$. For a nonzero target $y$, its norm fibre
is $\{(x,1_L):Nx=y\}$. These are exact set fibres; no norm
surjectivity on the original adeles has been assumed.

The norm is not an additive map when $n>1$. If some fibre of
places has at least two elements, put complementary zero-and-one
amplitudes in its two nonempty parts and zero amplitudes at all
other places. Each of the two norm adeles is zero, but the norm
of their sum has value one at this base place. The failure also
occurs on zero amplitudes with independent labels: at two places
above one $v$, use label patterns $(1_L,0_L)$ and $(0_L,1_L)$,
and use top labels at all other places above $v$. Their norm
labels are both bottom, whereas the norm label of their sum is
top. This argument works in every nontrivial $L$.

If there is only one place above the selected $v$, its local
degree is $n>1$. A separable field norm of degree $d>1$ on an
infinite field is not additive. Indeed additivity would give
$N(1+ta)=1+t^dN(a)$ for every scalar $t$; comparison with the
determinant polynomial forces the coefficient
$\operatorname{Tr}(a)$ of $t$ to vanish for every $a$.
The trace of a finite Galois extension is not the zero map.
Here is an elementary proof valid also when the characteristic
divides $d$. Distinct field automorphisms are linearly independent:
in a shortest nonzero relation $\sum_i c_i\sigma_i(x)=0$,
evaluate at $yx$ and subtract $\sigma_1(y)$ times the relation;
choose $y$ separating $\sigma_1$ and another automorphism.
This gives a shorter nonzero relation, a contradiction.
Applying this to the sum of all local automorphisms proves
that trace is nonzero. Since the local base field is infinite,
the polynomial comparison is legitimate, proving nonadditivity.

There is nonetheless an exact full distributive relation,
including every mixed term. For either carrier let $\iota$ and
$N$ denote the corresponding maps of GTL5. Semiring
distributivity gives
\begin{equation}\tag{GTL10}
 \iota N(X+Y)=\iota N(X)+\iota N(Y)
  +\sum_{\varnothing\ne A\subsetneq G}
       \left(\prod_{\gamma\in A}\gamma X\right)
       \left(\prod_{\gamma\notin A}\gamma Y\right).
\end{equation}
The mixed sum is Galois-fixed because the action permutes its
subsets, so it has an exact inverse image under $\iota$. This
defines the full norm-additivity defect in the original target,
with every mixed amplitude and label retained. Each subset term
in the local carrier has at $w$ the label
$\bigwedge_{\gamma\in A}\lambda_{\gamma^{-1}w}
 \wedge\bigwedge_{\gamma\notin A}\mu_{\gamma^{-1}w}$;
the sum joins these labels. Thus GTL10 retains the original
finite distributivity calculation rather than imposing additivity.

An actual additive receiving map is the trace, not the norm:
\begin{equation}\tag{GTL11}
 \begin{aligned}
 T_c(x,\lambda)&=(\operatorname{Tr}_{E/K}x,\lambda),\\
 (T_{\rm loc}(x,\boldsymbol\lambda))_v
  &=\left(\sum_{w\mid v}\operatorname{Tr}_{E_w/K_v}(x_w),
                              \bigvee_{w\mid v}\lambda_w\right).
 \end{aligned}
\end{equation}
Here the adelic trace uses the displayed local sum. We first
prove that both formulas land in their stated restricted carriers.
At every nonarchimedean place where $x_w\in\mathcal O_w$, all
its local Galois conjugates are integral. Their finite sum
$\operatorname{Tr}_{E_w/K_v}(x_w)$ is therefore integral and
belongs to $K_v$, hence belongs to $\mathcal O_v$. An adele has
only finitely many nonintegral finite coordinates; their images
under $\pi$ form a finite set. Outside that set every term of
the finite sum over $w\mid v$ belongs to $\mathcal O_v$, so the
output amplitude is an adele of $K$.
If its $v$-coordinate is nonzero, at least one input $x_w$ over
$v$ is nonzero. The constrained input carrier then forces
$\lambda_w=1_L$, and consequently
$\bigvee_{w\mid v}\lambda_w=1_L$ as required. If the output
coordinate is zero, its displayed join is an allowed zero label,
even when nonzero trace terms cancel. This proves landing for
$T_{\rm loc}$. For $T_c$, a nonzero output adele implies a
nonzero input adele and hence common label $1_L$; an output
zero adele retains the original permitted common label. Thus
the common carrier constraint is also preserved.
These maps are additive and linear over the corresponding base carrier
through $\iota$; finite lattice distributivity proves scalar
compatibility for $T_{\rm loc}$. Their comparisons are
$T_{\rm loc}\Delta_E=\Delta_KT_c$,
$T\iota(a,\lambda)=(na,\lambda)$, and
$\iota T(X)=\sum_{\gamma\in G}\gamma X$.
They are not asserted to replace the multiplicative norm.
The three pure label maps have the exact order adjunctions
\begin{equation}\tag{GTL12}
 J_\pi\dashv\pi^*\dashv M_\pi,
 \qquad
 J_\pi(\boldsymbol\lambda)_v=\bigvee_{w\mid v}\lambda_w,
 \quad \pi^*(\boldsymbol\mu)_w=\mu_{\pi(w)},
 \quad M_\pi(\boldsymbol\lambda)_v=\bigwedge_{w\mid v}\lambda_w.
\end{equation}
For example $\pi^*\mu\le\lambda$ means exactly that
$\mu_v\le\lambda_w$ for every $w\mid v$, or equivalently
$\mu\le M_\pi\lambda$; the join adjunction is the dual
inequality. In particular $M_\pi\pi^*=J_\pi\pi^*$ is the
identity. These are exact lattice maps even when norm is not
an additive morphism.

\subsection{Original local degrees, Haar modules and support localizations}

At nonarchimedean places retain additive Haar measures of mass
one on the rings of integers; at real and complex places use
their usual real additive Haar measures. For $z\in E_w^\times$
the multiplication map on the finite-dimensional $K_v$-vector
space has determinant $N_{E_w/K_v}(z)$. A coordinate Haar
measure differs from the selected local Haar measure by one
constant, which cancels in the scaling ratio. Consequently
\begin{equation}\tag{GTL13}
 |N_{E_w/K_v}(z)|_v=|z|_w,
 \qquad |a|_w=|a|_v^{d_w}\quad(a\in K_v^\times).
\end{equation}
These are the Haar-module absolute values: at a complex place
$|z|_w$ is the square of the usual complex modulus. For finite
places, multiplication by a uniformizer has module the inverse
residue-field cardinality, as follows by the finite index of its
image of the ring of integers. Thus GTL13 retains those finite
Haar factors, rather than using an unweighted complex or finite norm.
On ideles it gives the exact products
\begin{equation}\tag{GTL14}
 |N g|_K=\prod_v\prod_{w\mid v}|N_{E_w/K_v}(g_w)|_v
       =\prod_w|g_w|_w=|g|_E,
 \qquad |\iota c|_E=|c|_K^n.
\end{equation}
Only finitely many nonarchimedean factors differ from one.
These equalities concern ideles, where the products are defined;
no global absolute value is assigned here to an arbitrary adele.
Nor is the nonlinear field norm assigned an additive Haar
Jacobian by these multiplication-map identities.

At the first support localization the common and local carriers
become $L$ and $L^{\Sigma_F}$, as in PLP6. The common norm and
inclusion are identity on $L$; the local norm and inclusion become
$M_\pi$ and $\pi^*$ of GTL12. For a prime lattice ideal $J$ retain
the actual uniform-denominator localization
\begin{equation}\tag{GTL15}
 \mathscr L_{F,J}=(L^{\Sigma_F})
             [\operatorname{diag}(L\setminus J)^{-1}],\qquad
 [\lambda]=[\mu]\ \Longleftrightarrow\quad
 \exists\nu\notin J\ \forall v:\lambda_v\wedge\nu=\mu_v\wedge\nu.
\end{equation}
Finite meets over $w\mid v$ commute with this one uniform meet
by $\nu$, so $M_\pi$ descends to
$\mathscr L_{E,J}\to\mathscr L_{K,J}$. Pullback also descends,
and the two maps still compose to the identity on
$\mathscr L_{K,J}$. The join map descends by lattice
distributivity. Coordinate localization to $\prod_v L_J$
commutes with these finite-fibre operations, but that coordinate
map is not silently declared injective; PLP7--9's infinite-product
defect remains part of the comparison. Norm remains only a
multiplicative map at this localization when its meet map is
not join-preserving.

\subsection{Class groups, the right cocycle, and the correctly typed transfer}

Set $C_F=\mathbb A_F^\times/F^\times$. The inclusion of ideles
induces an injective map $i_C:C_K\to C_E$. Indeed if a base idele
becomes the principal idele of $a\in E^\times$, that principal
idele is Galois-fixed; hence $\gamma(a)=a$ for every $\gamma$ and
$a\in K^\times$. This proves injectivity without assuming an
injection of unrelated class-group quotients.

In fact $i_C$ identifies $C_K$ with $C_E^G$. Here is the exact
algebraic descent argument. For a fixed class with representative
$x\in\mathbb A_E^\times$, write
$c_\gamma=\gamma(x)/x\in E^\times$.
The cocycle law is
$c_{\gamma\delta}=c_\gamma\gamma(c_\delta)$.
Independence of the automorphisms, proved above, supplies an
$a\in E$ such that
$b=\sum_{\gamma\in G}c_\gamma^{-1}\gamma(a)\ne0$.
The cocycle law gives $\delta(b)=c_\delta b$, so $x/b$ is fixed.
Fixed ideles are exactly the diagonal base ideles; local fixed
fields and the integral-unit condition prove this at every place.
Thus the class is in $i_C(C_K)$. This is the multiplicative
Hilbert--90 calculation, included here to make the exact map explicit.

Let the relative Weil group be the original extension
\begin{equation}\tag{GTL16}
 1\longrightarrow C_E\longrightarrow W_{E,K}
             \longrightarrow G\longrightarrow1.
\end{equation}
Choose the source's section $s_\alpha$, with $s_1=1$, and keep
the right-hand cocycle and conjugation convention
\begin{equation}\tag{GTL17}
 a_{\alpha,\beta}=s_{\alpha\beta}^{-1}s_\alpha s_\beta\in C_E,
 \quad s_\alpha s_\beta=s_{\alpha\beta}a_{\alpha,\beta},
 \quad s_\alpha g s_\alpha^{-1}=\alpha(g),
 \quad x^\beta=\beta^{-1}(x).
\end{equation}
Associativity gives the full twisted cocycle law
\begin{equation}\tag{GTL18}
 a_{\alpha\beta,\gamma}(a_{\alpha,\beta})^\gamma
       =a_{\alpha,\beta\gamma}a_{\beta,\gamma},
 \qquad a_{1,\alpha}=a_{\alpha,1}=1.
\end{equation}
Indeed expand $(s_\alpha s_\beta)s_\gamma$ by moving its
$a_{\alpha,\beta}$ to the right of $s_\gamma$; this gives
$s_{\alpha\beta\gamma}a_{\alpha\beta,\gamma}
 (a_{\alpha,\beta})^\gamma$. Expanding the other bracketing
gives the other side. These are right-cocycle formulas; replacing
them by an untwisted product would change the group.

For precision define the ordinary finite-index group transfer
$\operatorname{Ver}:W_{E,K}\to C_E$ using these representatives.
For $w\in W_{E,K}$ write
$w s_\beta=s_{\pi_w(\beta)}h_{w,\beta}$, $h_{w,\beta}\in C_E$,
and put $\operatorname{Ver}(w)=\prod_\beta h_{w,\beta}$.
Multiplying two such expressions gives
$h_{wz,\beta}=h_{w,\pi_z(\beta)}h_{z,\beta}$, proving that
Ver is a group homomorphism because $C_E$ is abelian and
$\pi_z$ permutes the finite index set. Replacing representatives
by $s_\beta b_\beta$ changes each factor to
$b_{\pi_w(\beta)}^{-1}h_{w,\beta}b_\beta$; their product is
unchanged. Conjugating both $w$ and the representative system
therefore conjugates Ver's value. Since a homomorphism into an
abelian group is invariant under conjugating its argument, its
values are fixed under $G$. Thus Ver has image in $C_E^G$.

The correctly typed transfer in the source is
\begin{equation}\tag{GTL19}
 t=i_C^{-1}\operatorname{Ver}:W_{E,K}\longrightarrow C_K,
 \qquad
 i_C(t(s_\alpha))=\prod_{\beta\in G}a_{\alpha,\beta},
 \qquad t(g)=N(g)\quad(g\in C_E).
\end{equation}
The product of cocycle values in the middle expression is
literally in $C_E$, which is why $i_C$ is retained. For the last
equality, $g s_\beta=s_\beta\beta^{-1}(g)$ gives
$\operatorname{Ver}(g)=\prod_\beta\beta^{-1}(g)=i_C(Ng)$;
injectivity of $i_C$ then proves it. This also proves the
compatibility in GTL19 rather than trying to deduce it from
surjectivity alone. The deep class-field input, precisely as
stated in the cited CCM \texttt{transferhom} passage, is that
this $t$ is surjective. Denote $t_\alpha=t(s_\alpha)$. Its
homomorphism law and GTL17 give
\begin{equation}\tag{GTL20}
 t_\alpha t_\beta=t_{\alpha\beta}N(a_{\alpha,\beta}).
\end{equation}

\subsection{The full twisted monoid and its transfer}

For $F=K,E$ let
$X_{F,c}=S_{F,c}/F^\times$ and
$X_{F,\mathrm{loc}}=S_{F,\mathrm{loc}}/F^\times$,
where each principal nonzero element acts with top label.
These are orbit sets with well-defined multiplicative monoids:
changing representatives by nonzero principal factors does not
change a product orbit. No addition on these orbit sets is
assumed. The maps $N$, $\iota$, and $\Delta$ descend because
norm takes a principal element to its nonzero field norm and
inclusion takes principal elements to principal elements.
The descended inclusion is injective. For nonzero amplitudes,
an equality $\iota(b)=a\iota(c)$ with $a\in E^\times$ forces
$(\gamma(a)-a)\iota(c)=0$ for every $\gamma$. At one nonzero
coordinate this forces $\gamma(a)=a$, since a nonzero global
field element is nonzero in every completion. Thus $a\in K^\times$.
At zero amplitudes labels are unchanged by every principal action,
and injectivity follows directly from the label pullback.

The full fixed-orbit descent is
\begin{equation}\tag{GTL21a}
 \iota_c(X_{K,c})=(X_{E,c})^G,
 \qquad
 \iota_{\rm loc}(X_{K,\mathrm{loc}})
                         =(X_{E,\mathrm{loc}})^G.
\end{equation}
The inclusions already proved therefore identify the base orbit
monoids with these fixed orbit monoids. Here fixedness is for
the specified $G$ action on the full labelled orbit; it is not
only fixedness of the amplitude or of its zero pattern.

To prove GTL21a, the included base orbits are fixed by GTL6.
Conversely, choose a representative $X$ of a fixed orbit in
either $E$ carrier, and write $x$ for its amplitude adele.
First suppose $x\ne0$. For each $\gamma\in G$ there is a
$c_\gamma\in E^\times$ such that $\gamma X=c_\gamma X$.
It is unique: an equality $c x=d x$ at a nonzero coordinate
implies $c=d$, because $E$ embeds in that local field.
Associating $\gamma\delta X$ in two ways consequently gives
$c_{\gamma\delta}=c_\gamma\gamma(c_\delta)$.
The explicit Hilbert--90 calculation preceding GTL16 applies
to this cocycle, even though the adele $x$ need not be an idele.
In detail, independence of the distinct automorphisms gives
$a\in E$ with
$b=\sum_{\gamma\in G}c_\gamma^{-1}\gamma(a)\ne0$.
Since
$\delta(c_\gamma^{-1})
       =c_\delta c_{\delta\gamma}^{-1}$,
reindexing this finite sum gives $\delta(b)=c_\delta b$.
It follows that the representative $X/b$ is fixed by every
$\delta$. Multiplication by the principal element $b^{-1}$
preserves the adelic restricted condition and leaves every
support label unchanged: a nonzero global element is a local
unit at all but finitely many finite places and acts with top
label.

For completeness, the equality $\gamma X=c_\gamma X$ on the
place-labelled carrier gives
$\lambda_{\gamma^{-1}w}=\lambda_w$ at every place, including
the zero coordinates. Galois transitivity therefore makes
these labels constant on each fibre $w\mid v$. On the
common-label carrier the single label was unchanged from the
start. GTL6 now gives $X/b=\iota(Y)$ for a representative $Y$
in the corresponding $K$ carrier: fixed amplitudes descend
through the local fixed fields, and the descended labels obey
the original constraint. This exhibits the required base orbit.

If $x=0$, every principal nonzero element acts trivially on the
entire zero-labelled representative. Thus fixedness of its
orbit is actual fixedness of its labels. A common label descends
unchanged; a place-labelled tuple descends because it is
constant on each fibre of $\pi$. This covers all zero amplitudes
and all lattice labels without an extra choice of a cocycle.
Together with the preceding injectivity proof this establishes
both equalities of multiplicative monoids in GTL21a. No addition
or topological quotient identification is needed for this descent.

Use either one of the two $E$ orbit monoids and call it $X$.
The natural $C_E$ action is by top-supported idele multiplication;
the $G$ action is the one already specified. Define
\begin{equation}\tag{GTL21}
 (\alpha,x)\star(\beta,y)
       =(\alpha\beta,a_{\alpha,\beta}x^\beta y)
                     \quad\text{on }G\times X.
\end{equation}
Its two bracketings have the respective second coordinates
\[
 a_{\alpha\beta,\gamma}(a_{\alpha,\beta})^\gamma
         x^{\beta\gamma}y^\gamma z,
 \qquad
 a_{\alpha,\beta\gamma}a_{\beta,\gamma}
         x^{\beta\gamma}y^\gamma z.
\]
Here $(x^\beta)^\gamma=x^{\beta\gamma}$ by the original inverse
action convention. GTL18 proves associativity, and $(1,1)$ is
the unit. The original Weil group identifies with the submonoid
$G\times C_E$ via $(\alpha,g)\mapsto s_\alpha g$.

The left and right actions of the original group on this monoid are
therefore exactly
\begin{equation}\tag{GTL22}
 (s_\alpha g)(\beta,x)
   =(\alpha\beta,a_{\alpha,\beta}g^\beta x),\qquad
 (\alpha,x)(s_\beta g)
   =(\alpha\beta,a_{\alpha,\beta}x^\beta g).
\end{equation}
They commute by the proved associativity. Common labels are
retained by each group action. For local labels the left action
leaves the $x$ labels unchanged, whereas the right action with
$s_\beta g$ replaces the label at $w$ by the old label at
$\beta w$, exactly as $x^\beta=\beta^{-1}x$ requires.

Define the full supported transfer by
\begin{equation}\tag{GTL23}
 \mathcal T_c(\alpha,x)=t_\alpha N_c(x),\qquad
 \mathcal T_{\rm loc}(\alpha,x)=t_\alpha N_{\rm loc}(x).
\end{equation}
The factors $t_\alpha$ act through their idele classes with top
labels. These maps are well-defined and multiplicative for the
twisted monoids GTL21. In fact, applying one to a product gives
$t_{\alpha\beta}N(a_{\alpha,\beta})N(x^\beta)N(y)$;
norm is Galois-invariant and GTL20 makes this
$t_\alpha t_\beta N(x)N(y)$. All label products agree because
finite fibre meets commute with the place permutation.
Consequently
\begin{equation}\tag{GTL24}
 \mathcal T(wxz)=t(w)\mathcal T(x)t(z),\qquad
 \mathcal T_{\rm loc}(\alpha,\Delta_E x)
                          =\Delta_K\mathcal T_c(\alpha,x).
\end{equation}
This is the exact lifted equivariance and norm/diagonal square.
For a fixed $\alpha$, the map $x\mapsto t_\alpha N(x)$ is not
claimed to be a unital monoid homomorphism: its two product
expressions differ by one factor $t_\alpha$. Multiplicativity
in GTL23 refers to the full twisted monoid, not the untwisted
direct product or an additive quotient.

\subsection{Every one-zero-place label and the transferred classical orbits}

Let $a^{(w)}\in\mathbb A_E$ have amplitude zero at $w$ and one
at every other place. Its common-carrier lift has only top label.
In the separately labelled carrier write $a^{(w)}_\lambda$ for
the tuple with $(0,\lambda)$ at $w$ and $(1,1_L)$ elsewhere.
The exact norms are
\begin{equation}\tag{GTL25}
 N_c(a^{(w)},1_L)=(a^{(\pi(w))},1_L),\qquad
 N_{\rm loc}(a^{(w)}_\lambda)=a^{(\pi(w))}_\lambda.
\end{equation}
There is one zero local factor above $\pi(w)$; every other
factor is the norm of one. Thus its label meet is exactly
$\lambda$. The degree $d_w$ has not been replaced by a
multiplicity of zero labels: its repeated meet equals $\lambda$
by the idempotence of the original lattice operation.

Define the common and local classical orbit sets respectively by
\[
 \Xi_{F,c}=\bigcup_{v\in\Sigma_F}C_F[a^{(v)}_{1_L}],
 \qquad
 \Xi_{F,\rm loc}=\bigcup_{v\in\Sigma_F,\lambda\in L}
                                      C_F[a^{(v)}_\lambda].
\]
The common set embeds as the top-labelled part of the second.
Every label is preserved by idele multiplication. Formula GTL25
gives, for $g\in C_E$,
\begin{equation}\tag{GTL26}
 \mathcal T(\alpha,g a^{(w)}_\lambda)
       =t(s_\alpha g)a^{(\pi(w))}_\lambda.
\end{equation}
Here the common case has its required value $\lambda=1_L$.
For any $v$, choose a place $w\mid v$; for a target idele class
$c\in C_K$, the cited transfer surjectivity supplies
$s_\alpha g$ with $t(s_\alpha g)=c$. Hence
\begin{equation}\tag{GTL27}
 \mathcal T_c:G\times\Xi_{E,c}\twoheadrightarrow\Xi_{K,c},
 \qquad
 \mathcal T_{\rm loc}:G\times\Xi_{E,\rm loc}
                                  \twoheadrightarrow\Xi_{K,\rm loc}.
\end{equation}
These are the exact supported lifts of the source's classical
transfer surjection. They do not rely on surjectivity of the
norm $C_E\to C_K$, which is a different map.

The orbit fibres can also be specified without asserting an
unproved covering degree. Let $H_w$ be the image of
$E_w^\times$ in $C_E$ and $H_v$ the image of $K_v^\times$ in
$C_K$. These embeddings are injective: a principal idele equal
to one at a different place has principal field element one.
The stabilizer of $[a^{(w)}_\lambda]$ is $H_w$, since a
principal comparison forces its idele to equal that principal
element at every nonzero coordinate. Thus its orbit is $C_E/H_w$,
with the same statement over $K$. Since $N(H_w)\subset H_v$,
the exact fibre over $[c a^{(v)}_\lambda]$ is the set
\begin{equation}\tag{GTL28}
 \coprod_{w\mid v}\ \coprod_{\alpha\in G}
  \{gH_w:\ t_\alpha N(g)\in cH_v\},
\end{equation}
with label precisely $\lambda$. Different zero places and
different labels cannot be identified by principal or idele
multiplication. The formula is a set decomposition; it does
not replace the original ambient orbit topology by a
disjoint-union topology, or assert a local-homeomorphism property.

Finally let $C_{F,1}=\ker(|\cdot|_F)$. GTL14 gives
$N(C_{E,1})\subset C_{K,1}$. Right multiplication by $C_{E,1}$
in GTL22 changes only the $x$ coordinate, and GTL24 therefore
gives the descended surjections
\begin{equation}\tag{GTL29}
 G\times(\Xi_{E,c}/C_{E,1})\twoheadrightarrow
                      \Xi_{K,c}/C_{K,1},\qquad
 G\times(\Xi_{E,\rm loc}/C_{E,1})\twoheadrightarrow
                      \Xi_{K,\rm loc}/C_{K,1}.
\end{equation}
Every label remains unchanged in these quotients. For a single
zero-place local label at a support prime $J$, the equality
criterion is exactly $[\lambda]_J=[\mu]_J$, as in PLP10;
GTL25 transports that same criterion from $w$ to $\pi(w)$.
The original common point instead restricts to top in $L_J$.
These are consequences of the original arbitrary-lattice maps,
not a replacement of the support spectrum by a single point.

\paragraph{Precisely retained scope.}
Inclusion and diagonal comparison are unital semiring maps;
Galois actions are semiring automorphisms; norm is a unital
multiplicative map with exact mixed-additivity defect GTL10;
trace is the additive base-linear companion GTL11. On the orbit
sets only their specified multiplicative structures are used.
The full Weil cocycle and the typed class-group inclusion remain
in GTL17--24. Classical transfer surjectivity is the cited human
source input; the labelled fibres and surjections are its proved
receiving maps. No additive field-norm assertion, new norm
surjectivity, support-label Haar measure, or positivity theorem
is inferred.
